\chapter{Conclusioni}
\label{ch:conclusioni}

Questa tesi ha ricostruito diciassette anni di attività donazionale del DIMT
di Trieste e Gorizia --- 353\,671 donazioni di 43\,299 donatori, dal 2009
al 2025 --- per stimare quanto sangue il sistema
può attendersi di raccogliere, e da chi. Le due province coprono circa un
quinto del bacino donazionale del Friuli Venezia Giulia e non ne riproducono
le caratteristiche; le conclusioni che seguono vanno lette come relative al
territorio osservato.

\section*{Cosa si è trovato}

\textbf{Il sistema si contrae.} Il volume annuale
cala del $28{,}1\%$ fra il 2009 e il 2025, con la flessione concentrata nella
prima metà del periodo e un assestamento successivo. Rapportando le donazioni
alla popolazione residente di Trieste e Gorizia in età donabile --- una
scomposizione che questo lavoro introduce e che la sola serie assoluta non
consente --- si ottiene che il calo demografico spiega solo il $12{,}2\%$
della contrazione, mentre il restante $87{,}8\%$ è imputabile alla perdita di partecipazione. 

\textbf{Il fronte critico è contenuto in una fascia d'età.} A parità di popolazione
i 20--24enni sono la fascia con il tasso di donazione più alto in ogni anno
osservato, e costituiscono metà di tutti gli ingressi. Le fasce che crollano
sono le 35--44, che perdono fra il 36 e il 41 per cento del tasso di
partecipazione, mentre in corrispondenza delle fasce 55--64 si raggiunge uno stallo. 
Nel 2009 il profilo per età mostra un altopiano fra i 35 e i 49 anni,
che sparisce nel 2025. Il sistema, quindi,
non ha perso le nuove leve, ma i donatori di mezza età.

\textbf{Mantenere i giovani all'interno del sistema è difficile.} 
Il $40{,}6\%$ dei donatori che
entra fra i 18 e i 25 anni non dona più dopo il primo anno, e il
$52{,}3\%$ smette dopo il secondo. Il divario rispetto a chi entra più
tardi esiste ma è modesto --- rapporto di $1{,}22$ anziché $1{,}5$ --- una
volta rimossi dal confronto i donatori intercettati a metà carriera, che per
costruzione sono dei sopravvissuti. La funzione di rischio è decrescente:
superati i primi anni, il comportamento donazionale si consolida.

\textbf{Il margine sui donatori già attivi è vincolato.} I due effetti più
forti del modello di frequenza --- genere e tipologia di donazione --- sono
frutto di paletti regolamentari: intervalli minimi fra donazioni e
tetti annui differenziati per sesso. E controllando per gli anni di carriera, si
nota ciò che distingue un donatore assiduo è da
quanto tempo dona, indipendentemente dall'età anagrafica. 
Coerentemente, l'intensità individuale non cresce: le
donazioni per donatore attivo restano fra $1{,}84$ e $1{,}97$ per tutti e
diciassette gli anni.

\section*{Cosa significa operativamente}

I quattro punti convergono su un'indicazione unica. Il volume non si difende
chiedendo di più a chi già dona, perché quel margine è chiuso in alto da regole
cliniche e non da mancanza di disponibilità. Si difende agendo sul \emph{numero} di
donatori, e si può fare in due maniere:

\begin{enumerate}
    \item \textbf{secondo contatto}: quasi metà dei giovani che donano una
volta non torna, ed è qui dove si massimizza il rischio di abbandono . Un
richiamo rivolto a chi ha donato una volta e non si è ripresentato agisce sul
momento in cui la perdita è più probabile e in cui il valore recuperato ---
decenni di carriera potenziale --- è massimo.
    \item \textbf{rientro della fascia di mezza età}, che è il punto dove
la perdita si è effettivamente verificata. Chi ha smesso di donare a quarant'anni è,
a differenza di un neodonatore, una persona conosciuta dal sistema, 
che ha già superato la selezione clinica e che ha una storia donazionale registrata.
È il bacino su cui un'azione di richiamo ha il costo di contatto più basso.
\end{enumerate}

La proiezione al 2031 --- circa 97\,200 donazioni complessive nel sessennio,
con un primo anno più basso del 2025 e poi un calo di poco meno di un punto
percentuale l'anno --- va letta in questa
chiave: è la risposta alla domanda ``cosa succede se nulla cambia'', non una
previsione di ciò che accadrà. Il suo valore è di fissare la linea di base
rispetto a cui misurare un intervento.

\section*{Cosa non si è potuto stabilire}

La domanda aperta più rilevante riguarda le fasce 35--44: se il loro crollo sia
un effetto di \emph{età} o di \emph{coorte}. Se le generazioni oggi
quarantenni hanno sempre donato meno delle precedenti, il fenomeno è
generazionale e si sposterà con loro; se invece hanno donato e hanno smesso, è
recuperabile. Le due letture portano a politiche diverse, e i dati permettono
di porre la domanda ma non, attualmente, di trovarne una risposta: come discusso nella
\cref{sec:glm_spec}, età, periodo e coorte non sono identificabili
simultaneamente senza un'assunzione esterna. Il panel donatore-anno costruito
qui è lo strumento adatto ad affrontarla, ma va detto che i dati raccolti
finora non offrono nemmeno un indizio utilizzabile: nella variante del modello
che include la coorte di ingresso, i coefficienti \emph{crescono} verso le
generazioni più recenti, il che non depone per una minore propensione
generazionale --- riflette piuttosto il fatto che quelle coorti sono osservate
soltanto nei loro primi anni, che sono i più attivi. È la ragione per cui il
\cref{ch:glm} rifiuta di leggerlo come effetto generazionale.

\section*{Un'osservazione sul metodo}

Le conclusioni di questa tesi dipendono, in misura
sorprendente, da tre scelte facilmente classificabili come "tecniche":
quale delle due colonne disponibili identifichi il donatore; se un campo
amministrativo compilato in modo disomogeneo fra le sedi sia un requisito di
ammissione o una covariata; e su quale unità di osservazione si misuri un
conteggio.

Nessuna delle tre è una questione statistica quanto amministrativa,
ma tutte, in base alla scelta presa da chi svolge lo studio, hanno la
possibilità di cambiare i risultati finali. In particolare la seconda
aveva il potere di, essenzialmente, \textit{eliminare} un terzo delle
donazioni isontine --- perché la compilazione di quel campo è una prassi 
di sede, non un attributo del donatore.
